\slajdid{04-s036}
\begin{frame}[label=04-s036]
\gusttekst\setlength{\abovedisplayskip}{4pt}\setlength{\belowdisplayskip}{4pt}
UOČITE

Za računanje negativnih vrednosti u drugom komplementu binarnog brojnog sistema

Znamo da je prikaz broja negativan

Pravu vrednost dobijamo komplementiranjem i stavljanjem znaka minus ispred
\[\begin{aligned}
&-\left(2^n-b_{n-1}b_{n-2}\,b_{n-3}\ldots b_1\,b_0\right)\\
&=-\left(2^n-b_{n-1}2^{n-1}-b_{n-2}\,b_{n-3}\ldots b_1\,b_0\right)\\
&=-\left((2-b_{n-1})2^{n-1}-b_{n-2}\,b_{n-3}\ldots b_1\,b_0\right)\\
&=-(2-b_{n-1})2^{n-1}+b_{n-2}b_{n-3}\ldots b_1\,b_0\\
&=(b_{n-1}-2)2^{n-1}+b_{n-2}b_{n-3}\ldots b_1\,b_0
\end{aligned}\]
\[b_{n-1}=1\]
\[-b_{n-1}\,2^{n-1}+b_{n-2}b_{n-3}\ldots b_1\,b_0\]
PRELAZAK\qquad $\textcolor{DEcrvena}{(-b_{n-1})\,b_{n-2}b_{n-3}\ldots b_1\,b_0}$
\end{frame}
